% =====================================================================
% Table 2 --- the closure conditions, in simplified formal form
% ---------------------------------------------------------------------
% What Section 3's closing sentence promises: Table 1 recast "in terms of
% specific constellations of closure". Here the content is the CLOSURE
% CONDITIONS themselves -- two dimensions, five closures -- with the equation
% each one pins and the margin that is left to adjust.
% Alternative shape (if you prefer the 3x2 grid kept): the same conditions in
% the row/column headers of Table 1, with the qualitative outcomes in the
% cells. Say the word and I swap it.
% =====================================================================
\begin{table}[htbp]
\centering
\small
\begin{tabular}{@{}l l l@{}}
\toprule
\textbf{Closure} & \textbf{Condition (simplified)} & \textbf{What adjusts} \\
\midrule
\multicolumn{3}{@{}l}{\textit{Short run --- what the labour market pins}} \\[2pt]
factor supply fixed
  & $\sum_i L_i = \bar L$
  & wages and prices \\[3pt]
real wage fixed
  & $w/P = \bar w$
  & employment $L$ \\[5pt]
\multicolumn{3}{@{}l}{\textit{Money --- how the programme is financed}} \\[2pt]
repurposing of consumption (F1)
  & $\sum_i p_i c_i = E_h$
  & the composition of $c$ \\[3pt]
refinancing via taxes (F2)
  & $\sum_i p_i g_i = T$
  & the tax bill $T$ \\[3pt]
foreign debt (F3)
  & $\sum_i p_i g_i = F$
  & the external transfer $F$ \\
\bottomrule
\end{tabular}
\caption{Table~\ref{tab:overview} recast as specific constellations of
closure: the conditions each closure pins, in simplified form. The full
matrix follows in Section~\ref{sec:application}.}
\label{tab:closure-conditions}
\end{table}
